\begin{abstract}
Although large-scale unsupervised language models (LMs) acquire extensive world knowledge and develop certain reasoning abilities, exerting fine-grained control over their outputs remains challenging owing to the entirely unsupervised paradigm in which they are trained. To address this, prior approaches typically collect human-provided preference data indicating which model outputs are superior, subsequently fine-tuning the base LM to conform to these preferences—most frequently utilizing reinforcement learning from human feedback (RLHF). Nonetheless, the RLHF process involves considerable complexity and instability: it first necessitates training a reward model that approximates human judgments and then updates the unsupervised LM through reinforcement learning so as to maximize this proxy reward, all while constraining deviation from the initial model. In this work, we propose a novel parameterization of the RLHF reward model that permits closed-form extraction of the associated optimal policy. This reformulation enables solving the standard RLHF problem using only a straightforward classification loss. The introduced algorithm, termed \textit{Direct Preference Optimization} (DPO), is robust, efficient, and requires minimal computation, obviating the need for sampling from the LM during fine-tuning or extensive hyperparameter searches. Empirical results demonstrate that DPO effectively aligns LMs with human preferences, achieving parity with or superiority to previous methods. In particular, DPO fine-tuning surpasses PPO-based RLHF in sentiment control and matches or enhances output quality in summarization and single-turn dialogue tasks, all with greater ease of use and training efficiency.
\end{abstract}

\section{Introduction}
Large unsupervised language models (LMs), when trained on massive datasets, have demonstrated remarkable and unexpected abilities~\citep{chowdhery2022palm, brown2020language, touvron2023llama, bubeck2023sparks}. Nonetheless, because these models are trained using data created by humans with divergent intentions, competencies, and motivations, they inevitably absorb a mixture of desirable and undesirable traits. For instance, in the context of AI-assisted programming, it is beneficial for a model to \textit{recognize} prevalent coding errors to correct them; however, we would still prefer that the code the model produces reflects the superior—albeit infrequently represented—programming practices embedded within the training set. In a similar vein, even though it can be valuable for a language model to possess knowledge of misconceptions widely held by the public (such as one believed by half of people), we do not want the model to propagate that misinformation in half of the responses it provides on the topic. Therefore, it is essential to delineate the model's \emph{intended outputs and behaviors} from the broad spectrum of its underlying \textit{knowledge and capabilities} to ensure the creation of AI systems that are reliable, high-performing, and amenable to human direction~\citep{ouyang2022training}. Although prevailing techniques typically align LMs with human preferences through reinforcement learning (RL), we demonstrate that the RL-based objective frequently employed can, in fact, be equivalently optimized via a straightforward binary cross-entropy objective, thus streamlining the entire preference learning process.

In summary, dominant approaches employ collections of carefully selected human preference data to imbue language models with attributes that humans value, especially those contributing to safety and usefulness. This phase of preference learning takes place subsequent to an initial period of large-scale unsupervised pre-training on vast text corpora. While a direct route to aligning models with human preferences involves supervised fine-tuning using examples of desirable responses provided by humans, reinforcement learning from human (or AI) feedback (RLHF/RLAIF; \citep{christiano2017deep, bai2022constitutional}) has become the most effective paradigm. Within RLHF, a reward model is learned based on annotated human preferences, and this model guides RL optimization of the language model, encouraging the generation of high-reward outputs while discouraging significant divergence from the base model. Despite the considerable advances produced by RLHF, particularly in dialogue and programming domains, this framework introduces notable complexity relative to supervised learning by necessitating the training of several language models and ongoing policy sampling throughout training, which significantly increases the computational overhead.

In this work, we demonstrate a method for optimizing language models to reflect human preferences directly, bypassing the need for explicit reward modeling or reinforcement learning techniques. We introduce \textit{{Direct Preference Optimization} (DPO)}, a novel algorithm that implicitly targets the same objective as standard RLHF approaches—reward maximization under a KL-divergence constraint—while maintaining a streamlined implementation and training procedure. At a high level, DPO operates by boosting the log probability ratio of responses favored by humans over those that are not, augmented by an adaptive, example-specific importance weighting mechanism designed to avoid the model deterioration observed when relying on a naive probability ratio approach. Similar to previous algorithms, DPO is grounded in a theoretical framework for preference modeling, such as the Bradley-Terry model \cite{bradley1952rankanalysis}, which quantifies the correspondence between a reward function and observed preference data. Unlike prevailing methods that utilize this model to construct a preference loss for reward model training, and subsequently optimize a policy with respect to the learned reward model, DPO reformulates the preference loss directly as a function of the policy through a change of variables. This allows DPO to optimize a policy based on binary cross entropy over datasets containing human-labeled preferences, resulting in the optimal policy with respect to an implicit reward model tailored to the data.

Our principal contribution is Direct Preference Optimization (DPO), an RL-free, straightforward approach for leveraging preference data to train language models. Empirical evaluations indicate that DPO matches or surpasses the efficacy of established methods—including PPO-based RLHF—for various preference-driven tasks such as sentiment control, summarization, and dialogue, scaling up to models with 6B parameters.

\section{Related Work}

As self-supervised language models increase in size, they demonstrate the capability to perform certain tasks without explicit task-specific training—either in a zero-shot setting \citep{radford2019language} or by utilizing few-shot prompts \citep{gpt3,megatron,chowdhery2022palm}. Nevertheless, their results on various downstream tasks and their capacity to conform to user intent are greatly enhanced when subjected to fine-tuning on corpora containing explicit instructions and completions written by humans \citep{mishra-etal-2022-cross,sanh2022multitask,chung2022scaling,thoppilan2022lamda}. This process, known as `instruction-tuning,’ equips LLMs with the ability to extrapolate to novel instructions not seen during training, and typically results in improved practicality and user satisfaction \citep{chung2022scaling}. Notably, even though instruction tuning has proven effective, collecting \textit{relative} assessments of response quality from humans is generally more feasible than amassing datasets of expert demonstrations. As a result, later studies have leveraged datasets constructed from human preference judgments to further fine-tune LLMs, leading to notable advances in tasks such as translation \citep{kreutzer-etal-2018-reliability}, summarization \citep{stiennon2022learning,ziegler2020finetuning}, story generation \citep{ziegler2020finetuning}, and following user instructions \citep{ouyang2022training,ramamurthy2023is}. This approach first involves training a reward model—typically a neural network—optimized for agreement with human preference data via models such as the Bradley-Terry framework \citep{bradley1952rankanalysis}. Subsequently, the language model undergoes reinforcement learning, aiming to maximize this reward using algorithms such as REINFORCE \citep{williams1992reinforce}, proximal policy optimization (PPO; \cite{schulman2017proximal}), or related techniques \citep{ramamurthy2023is}. In parallel, another research trajectory has emerged wherein LLMs, already refined for instruction following using human feedback, are further utilized to generate supplementary synthetic preference annotations targeting particular features like safety or benign behavior \citep{bai2022constitutional}. This is accomplished with minimal human input, typically limited to written guidelines informing the LLM’s rubric for annotation. Collectively, these strategies integrate two major research avenues: one focusing on reinforcement learning for diverse language modeling tasks~\citep{Ranzato2015SequenceLT,paulus2018a,wu2018learning}, and another centered on general principles for preference-based learning from humans \citep{christiano2017deep,kupcsik2018learning}. While using relative human preferences offers significant advantages, employing reinforcement learning to fine-tune large language models continues to present substantial practical obstacles; the present work introduces a theoretically-motivated methodology for optimizing relative preferences that bypasses the need for reinforcement learning.

Research on learning policies from preferences outside linguistic domains has been conducted within both bandit and reinforcement learning paradigms, with multiple methodologies emerging. One such framework is contextual dueling bandits (CDB; \cite{yue2012karmed,dudik2015contextual}), which relies on preferences or rankings over actions, as opposed to direct reward signals. Without absolute reward feedback, the theoretical framework of CDBs replaces the standard optimal policy definition with the concept of a \textit{von Neumann winner}: a policy that secures an expected win rate of at least 50\% against all alternative policies \citep{dudik2015contextual}. Notably, CDB methods acquire preference information sequentially in an online fashion, whereas policy learning from human feedback more commonly relies on a static, pre-collected dataset containing preference annotations for action pairs \citep{yan2022human}. 

In a similar vein, \textit{preference-based RL} (PbRL) derives training signals from binary preferences yielded by an \textit{unknown} `scoring' or utility function, in lieu of explicit rewards \citep{BusaFekete2014,ruiz2023dueling}. PbRL literature offers various strategies, including those designed for reusing preference data collected off-policy; however, these generally involve a two-step process: estimating the underlying scoring function (serving as a reward proxy), followed by its optimization \citep{jain2013learning,BusaFekete2014,christiano2017deep,sadigh2017active,kupcsik2018learning}. In contrast, we introduce a unified, single-stage procedure that directly updates the policy to satisfy observed preferences.

\section{Preliminaries}\label{section:prelims}

We provide a summary of the RLHF pipeline as described in \citeauthor{ziegler2020finetuning} and subsequent works (\citep{stiennon2022learning, bai2022training, ouyang2022training}), which typically encompasses three sequential stages: (1) supervised fine-tuning (SFT); (2) preference elicitation and reward modeling; and (3) reinforcement learning-based optimization.

\textbf{SFT}: The RLHF protocol generally starts by further training a pre-trained language model (LM) via supervised learning on a curated set of high-quality examples for the tasks of interest (such as summarization or dialogue), resulting in an adapted model, denoted $\pi^\text{SFT}$.

\textbf{Reward Modeling Phase}: During the subsequent stage, the SFT model is conditioned on various prompts $x$ to generate response pairs $(y_1, y_2)\sim \pi^\text{SFT}(y \mid x)$. These candidate pairs are assessed by human annotators, who express their preference between them, which is encoded as $y_w\succ y_l \mid x$, where $y_w$ and $y_l$ are, respectively, the more and less favored completions within the pair $(y_1, y_2)$. Although the true reward model $r^*(y, x)$ that informs these judgments remains unobserved, preference modeling can proceed using statistical frameworks such as the Bradley-Terry (BT) model \cite{bradley1952rankanalysis}, a widely used method. More general models for ranked data, like the Plackett-Luce model \citep{plackett1975analysis, luce2012individual}, are also applicable if more extensive ranking information is available. The BT model describes the distribution of human preferences $p^*$ as follows:

Given a static collection of preference comparisons $\mathcal{D}=\bigl\{x^{(i)}, y_w^{(i)}, y_l^{(i)}\bigr\}_{i=1}^N$ drawn according to $p^*$, we introduce a reward estimator $r_{\phi}(x, y)$ and fit its parameters using maximum likelihood estimation. This preference modeling can be reframed as a binary classification problem, yielding the negative log-likelihood loss function:

\begin{equation}\label{eq:reward_model}
    \mathcal{L}_R(r_{\phi}, \mathcal{D}) = -\mathbb{E}_{(x, y_w, y_l)\sim \mathcal{D}}\left[\log \sigma(r_{\phi}(x, y_w)- r_{\phi}(x, y_l))\right]
\end{equation}

where $\sigma$ denotes the logistic sigmoid. In the case of language models, it is typical to initialize $r_{\phi}(x, y)$ with parameters inherited from the SFT model $\pi^\text{SFT}(y \mid x)$ and to append a linear output head atop the final transformer representation, producing a scalar score interpreted as the reward~\cite{ziegler2020finetuning}. To achieve a reward function with reduced variance, it is customary to normalize the outputs such that $\mathbb{E}_{x,y\sim \mathcal{D}}[r_\phi(x, y)] = 0$ for any context $x$.

\textbf{RL Fine-Tuning Phase}: During fine-tuning via reinforcement learning, the trained reward model delivers the supervisory signal to guide the policy updates. Drawing on previous approaches~\citep{jaques2017sequence, jaques2020human}, the optimization objective takes the following form:

\begin{equation}\label{eq:RL}
\max_{\pi_{\theta}}  \mathbb{E}_{x\sim \mathcal{D}, y\sim \pi_{\theta}(y \mid x)}\left[r_{\phi}(x, y)\right] - \beta\mathbb{D}_{\textrm{KL}}\bigl[\pi_{\theta}(y\mid x)\mid \mid \pi_\text{ref}(y\mid x)\bigr],
\end{equation}

with $\beta$ modulating the penalty for diverging from a reference distribution $\pi_\text{ref}$, generally set as the original SFT policy $\pi^\text{SFT}$. The policy $\pi_\theta$ for the language model is correspondingly initialized from $\pi^\text{SFT}$. This KL regularization serves to anchor the policy close to distributions where the reward model's predictions are trustworthy, in addition to safeguarding generative variety and discouraging mode collapse towards only high-rewarded responses. As text generation is inherently discrete, the objective cannot be differentiated directly and thus requires reinforcement learning for its optimization. The conventional method~\cite{ziegler2020finetuning, stiennon2022learning, bai2022training, ouyang2022training} is to combine the reward and the KL penalty into a single function, ${r(x, y) = r_{\phi}(x, y) -\beta (\log \pi_{\theta}(y\mid x) - \log \pi_\text{ref}(y\mid x))}$, and apply PPO~\cite{schulman2017proximal} for maximization.

\section{Direct Preference Optimization}\label{sec:DPO}

Recognizing the practical obstacles associated with applying reinforcement learning to the large-scale fine-tuning of language models, we aim to devise a streamlined methodology that allows direct optimization over preferences. Unlike standard RLHF pipelines, which construct and subsequently maximize a reward model using RL algorithms, our procedure exploits a specific reward parameterization that admits an analytical solution for the optimal policy—bypassing the need for an explicit RL loop. Central to our approach is the utilization of a closed-form correspondence between reward functions and their induced optimal policies. This enables re-expressing the loss—originally formulated over rewards—as a policy-level loss. Adopting this change-of-variable principle obviates the requirement to train a standalone reward predictor, while retaining compatibility with established preference models like the Bradley-Terry model. Effectively, our architecture combines both the policy (language model) and an implicit reward within the same network.

\textbf{Deriving the DPO objective.} We begin with the reinforcement learning (RL) objective presented in Eq.~\ref{eq:RL}, considering an arbitrary reward function $r$, as is customary in previous literature. Building upon methods established in~\citep{peters2007reinforcement, peng2019advantage, korbak2022reinforcement, go2023aligning}, it is straightforward to confirm that the solution to the KL-regularized maximization of expected reward described by Eq.~\ref{eq:RL} admits the following form:

\begin{equation}\label{eq:op_policy}
    \pi_r(y\mid x) = \frac{1}{Z(x)}\pi_\text{ref}(y\mid x)\exp\left(\frac{1}{\beta}r(x, y)\right),
\end{equation}

with $Z(x) =\sum_{y}\pi_\text{ref}(y\mid x)\exp\left(\frac{1}{\beta}r(x, y)\right)$ denoting the partition function. A complete derivation is provided in Appendix~\ref{app:derivation1}. Directly estimating $Z(x)$ remains computationally costly, even when one substitutes the maximum likelihood estimate $r_\phi$ for the true reward $r^*$, as indicated by \citep{korbak2022reinforcement, go2023aligning}. This renders the use of Eq.~\ref{eq:op_policy} challenging for practical applications. Nevertheless, it is possible to rearrange Eq.~\ref{eq:op_policy} so that the reward can be rewritten in terms of its optimal policy $\pi_r$, reference policy $\pi_\text{ref}$, and the (still intractable) $Z(\cdot)$. Specifically, by taking the logarithm of both sides of Eq.~\ref{eq:op_policy} and performing straightforward algebraic manipulations, we obtain:

\begin{equation}\label{eq:main_eq}
    r(x,y) =\beta \log \frac{\pi_r(y\mid x)}{\pi_\text{ref}(y\mid x)} + \beta \log Z(x).
\end{equation}

We then apply this reformulation to the true reward $r^*$ and the associated optimal policy $\pi^*$. Importantly, the Bradley-Terry preference model depends solely on the reward difference between two candidate responses: ${p^*(y_1 \succ y_2 \mid x) = \sigma(r^*(x, y_1) - r^*(x, y_2))}$. Plugging in the reparametrized form of $r^*$ from Eq.~\ref{eq:main_eq} into the Bradley-Terry equation (Eq.~\ref{eq:bradley-terry}) causes the partition function terms to cancel, allowing us to express preference probabilities entirely in terms of $\pi^*$ and $\pi_\text{ref}$. Consequently, under the Bradley-Terry model, the optimal RLHF policy $\pi^*$ obeys the following relationship:

\begin{equation}\label{eq:objective}
    p^*(y_1\succ y_2 \mid x)=\frac{1}{1 + \exp\left(\beta \log \frac{\pi^*(y_2\mid x)}{\pi_\text{ref}(y_2\mid x)} - \beta \log \frac{\pi^*(y_1\mid x)}{\pi_\text{ref}(y_1\mid x)}\right)}
\end{equation}

Details of this derivation are found in Appendix~\ref{app:derivation2}. Although Eq.~\ref{eq:objective} relies on the Bradley-Terry model, analogous results can be deduced for broader preference models, such as the Plackett-Luce family~\citep{plackett1975analysis, luce2012individual}, as explained in Appendix~\ref{app:plackett_luce_models}.

Expressing the likelihood of observed human preference data as a function of the optimal policy (rather than the reward function) permits direct construction of a maximum likelihood objective for a learnable policy $\pi_\theta$. Mirroring the approach of reward modeling (cf. Eq.~\ref{eq:reward_model}), this leads us to the following policy-based objective:

\begin{equation}\label{eq:optimum_model}
    \mathcal{L}_\text{DPO}(\pi_{\theta}; \pi_\text{ref}) = -\mathbb{E}_{(x, y_w, y_l)\sim \mathcal{D}}\left[\log \

In this framework, we approximate an implicit reward function via an alternative parameterization, where the resulting optimal policy coincides with $\pi_\theta$. Furthermore, since our approach corresponds to training a reparametrized Bradley-Terry model, it inherits several theoretical guarantees, such as consistency provided that the preference data distribution satisfies appropriate conditions \cite{bong2022generalized}. Additional discussion of the theoretical underpinnings of DPO and its relationship to existing literature is presented in Section~\ref{sec:theory}.

\textbf{Mechanics of the DPO update.} To gain insight into DPO, it is instructive to examine the gradient of its objective function $\mathcal{L}_\text{DPO}$. The derivative with respect to the model parameters $\theta$ can be expressed as:

\begin{multline*}\label{eq:gradient}
    \nabla_\theta \mathcal{L}_\text{DPO}(\pi_\theta;\pi_\text{ref}) = \\ -\beta\mathbb{E}_{(x, y_w, y_l) \sim \mathcal{D}} \left[\underbrace{\sigma(\hat{r}_\theta(x, y_l) - \hat{r}_\theta(x, y_w))}_\text{assigns larger weight when reward model is incorrect} \left[\underbrace{\nabla_\theta\log \pi(y_w \mid x)}_\text{encourage $y_w$} - \underbrace{\nabla_\theta\log\pi(y_l \mid x)}_\text{discourage $y_l$}\right]\right],
\end{multline*}

where $\hat{r}_\theta(x, y) = \beta \log \frac{\pi_\theta(y \mid x)}{\pi_\text{ref}(y \mid x)}$ denotes the reward that is implicitly specified by the parameterized model $\pi_\theta$ and the reference $\pi_\text{ref}$ (for further discussion, see Section~\ref{sec:theory}). In essence, the loss gradient for $\mathcal{L}_\text{DPO}$ encourages higher probability for the preferred responses $y_w$ while reducing probability for less preferred alternatives $y_l$. Notably, the degree to which each instance affects the update is determined by how much the implicit reward model $\hat{r}_\theta$ erroneously favors the less preferred option, scaled by $\beta$. This mechanism reflects the extent to which the implicit reward fails to correctly rank the options, as well as the influence of the KL divergence constraint. Our empirical studies highlight the necessity of this weighting; omitting the coefficient leads to a degenerate language model (see Appendix Table~\ref{tab:unlikelihood_generations}).

\textbf{DPO overview.} The typical DPO workflow consists of the following steps: 1) For each prompt $x$, generate candidate completions $y_1, y_2 \sim \pi_\text{ref}(\cdot \mid x)$ and assign human preference labels to form an offline preferences dataset $\mathcal{D} = \{x^{(i)}, y_w^{(i)}, y_l^{(i)}\}_{i=1}^N$; 2) train the language model $\pi_\theta$ by minimizing the DPO loss $\mathcal{L}_\text{DPO}$ with respect to the fixed $\pi_\text{ref}$, the dataset $\mathcal{D}$, and a chosen $\beta$ parameter.

In actual applications, it is preferable to utilize existing publicly available preference datasets rather than produce fresh samples and solicit human evaluations. Since such datasets are often derived using $\pi^\text{SFT}$, we set $\pi_\text{ref} = \pi^\text{SFT}$ whenever possible. In cases where $\pi^\text{SFT}$ is not accessible, we instead fit $\pi_\text{ref}$ by maximizing the likelihood over the preferred completions $(x, y_w)$ according to $\mathcal{D}$, formally: ${\pi_\text{ref} = \argmax_{\pi}\mathbb{E}_{x, y_w \sim \mathcal{D}}\left[\log \pi(y_w \mid x)\right]}$. This method serves to reduce the discrepancy between the unavailable true reference distribution and the $\pi_\text{ref}$ employed in DPO. Further implementation details and hyperparameter choices are provided in Appendix~\ref{app:implementation}.

\section{Theoretical Analysis of DPO} In this section, we further interpret the DPO approach, present supporting theoretical results, and clarify how DPO addresses difficulties inherent to actor-critic frameworks used in RLHF (for example, PPO~\cite{schulman2017proximal}).

\label{sec:theory}

\subsection{Your Language Model Is Secretly a Reward Model} DPO circumvents the need for explicitly fitting a reward model or running a separate RL optimization by relying on a unified maximum likelihood formulation. The loss function in Eq. \ref{eq:main_eq} corresponds to a Bradley-Terry model where the rewards are given by $r^*(x, y) = \beta \log\frac{\pi^*_\theta(y \mid x)}{\pi_\text{ref}(y \mid x)}$. Training the parameterized model $\pi_\theta$ under this objective is equivalent, via a change of variables, to maximizing the likelihood of the reward model as in Eq. \ref{eq:reward_model}. In the subsections that follow, we develop the theoretical framework underlying this reparameterization, demonstrate that it does not restrict the expressiveness of the reward functions that can be learned, and prove that this enables the precise recovery of the optimal policy. To set the stage, we first introduce an equivalence relation over reward functions.

\begin{definition}
Two reward functions $r(x, y)$ and $r'(x, y)$ are defined as equivalent if and only if ${r(x, y)-r'(x, y) = f(x)}$ for some function $f$.
\end{definition}

It is straightforward to verify that this relation constitutes an equivalence relation, which subsequently divides the space of reward functions into equivalence classes. This leads us to the following two lemmas:

\begin{lemma}\label{lemma:same_prefrence}
Within the Plackett-Luce, and specifically the Bradley-Terry, preference models, any pair of reward functions from the same equivalence class result in identical preference distributions.
\end{lemma}

\begin{lemma}\label{lemma:same_policy}
    Any two reward functions belonging to a common equivalence class lead to the same optimal policy when considering the constrained RL setting.
\end{lemma}

The arguments for these results are elementary and can be found in Appendix \ref{app:lemma1}. The first lemma highlights the classic issue of under-specification within the Plackett-Luce model family \cite{plackett1975analysis}. As a consequence, it is common practice to enforce additional identifiability conditions to guarantee that the MLE estimates in Eq. \ref{eq:reward_model} are meaningful \cite{bong2022generalized}. The second lemma ensures that the optimal policy is invariant to the choice of reward function within an equivalence class; therefore, our principal objective becomes finding any representative from the optimal reward class. The following theorem is established in Appendix~\ref{app:thm1}:

\begin{theorem}\label{thm:main}
    Provided certain mild conditions are met, every reward class compatible with the Plackett-Luce (and specifically the Bradley-Terry) model can be expressed through the parameterization ${r(x, y) = \beta \log \frac{\pi(y\mid x)}{\pi_\text{ref}(y\mid x)}}$, where $\pi(y\mid x)$ is a model and $\pi_\text{ref}(y \mid x)$ is a fixed reference model.
\end{theorem}

\begin{sproof}
    Let us take any reward function $r(x, y)$ and its corresponding optimal model $\pi_r(y \mid x)$, as defined by Eq. \ref{eq:op_policy}. We demonstrate that some function within the equivalence class of $r$ can be written via the reparameterization above. Define the mapping $f$ as
\begin{equation}
    f(r; \pi_\text{ref}, \beta)(x, y) = r(x, y) - \beta\log\sum_{y}\pi_\text{ref}(y\mid x)\exp\left(\frac{1}{\beta}r(x, y)\right)
\end{equation}
This operator $f$ performs normalization by subtracting the log-partition function of $\pi_r$ from the reward. The subtracted normalization only depends on $x$, so $f(r; \pi_\text{ref}, \beta)(x, y)$ remains in the same equivalence class as $r(x, y)$. Substituting the form of $r$ provided by the right-hand side of Eq.~\ref{eq:main_eq} (which is generally applicable), yields $f(r; \pi_\text{ref}, \beta)(x, y) = \beta \log \frac{\pi_r(y\mid x)}{\pi_\text{ref}(y\mid x)}$. Hence, $f$ generates a representative in $r$'s equivalence class with the requisite structure, implying that the proposed reparameterization suffices to cover all relevant reward functions in this context.
\end{sproof}

An alternative interpretation of Theorem~\ref{thm:main} is that it uniquely identifies, within each equivalence class, the specific reward function selected by the DPO reparameterization. This is the reward function that satisfies:

\begin{equation}\label{eq:lag_p}
     \sum_{y}\underbrace{\pi_\text{ref}(y\mid x)\exp\left(\frac{1}{\beta}r(x, y)\right)}_{=\pi(y\mid x)\text{, using Thm.~\ref{thm:main} reparam.}} = 1,
\end{equation}

meaning that $\pi(y\mid x)$ is indeed a normalized probability distribution, assigning non-negative values summing to one. Inspecting Eq.~\ref{eq:lag_p} in light of Eq.~\ref{eq:op_policy}, we observe that it corresponds to the partition function for the optimal policy determined by the reward function $r(x, y)$. The central observation in the DPO algorithm is that by introducing specific constraints into the Plackett-Luce family (and, specifically, the Bradley-Terry model) of preference-based models—which are otherwise underdetermined—we can maintain the diversity of possible reward representations, yet guarantee that the optimal policy of Eq.~\ref{eq:op_policy} remains analytically manageable for all choices of $x$.

\subsection{Instability of Actor-Critic Algorithms} Our framework further enables an analysis of instability issues that occur with conventional actor-critic algorithms, such as PPO, which are commonly applied in RLHF. We adhere to the typical RLHF methodology and restrict attention to the reinforcement learning fine-tuning phase described in Section~\ref{section:prelims}. Here, we can relate our approach to the control as inference paradigm~\cite{levine2018reinforcement}, as applied to the constrained RL problem in~\ref{eq:RL}. Assuming a parameterized policy $\pi_{\theta}(y\mid x)$, we seek to minimize $\mathbb{D}_{\text{KL}}[\pi_{\theta}(y|x) \mid \mid \pi^*(y\mid x)]$, where $\pi^*$, determined by Eq.~\ref{eq:optimum_model}, is the optimal policy with respect to the reward $r_{\phi}(y, x)$. Straightforward algebraic manipulation yields the following maximization problem:

\begin{equation}\label{eq:AC}
    \max_{\pi_{\theta}}\mathbb{E}_{\pi_{\theta}(y\mid x)}\bigg[\underbrace{r_{\phi}(x, y) -\beta\log\sum_{y}\pi_\text{ref}(y\mid x)\exp\left(\frac{1}{\beta}r_{\phi}(x, y)\right)}_{f(r_{\phi}, \pi_\text{ref}, \beta)} - \underbrace{\beta\log\frac{\pi_{\theta}(y\mid x)}{\pi_\text{ref}(y\mid x)}}_{\text{KL}}\bigg]
\end{equation}

The objective considered here is identical to that in earlier studies 
\citep{ziegler2020finetuning, stiennon2022learning, bai2022training, ouyang2022training}, where optimization is performed using the DPO-equivalent reward from the reward class $r_{\phi}$. Within this framework, the normalization factor appearing in $f(r_{\phi}, \pi_\text{ref}, \beta)$ can be viewed as representing the soft value function for the reference policy $\pi_\text{ref}$. Although this normalization term is not necessary for characterizing the optimal solution, omitting it may result in a high-variance policy gradient, potentially leading to unstable training dynamics. To manage this normalization, one option is to approximate it with a learned value function, though such an approach presents its own optimization challenges. Another method employed in earlier works is to normalize the rewards relative to a human completion baseline, which effectively corresponds to a single-sample Monte-Carlo estimation of the normalization constant. In contrast, the reparameterization underlying DPO produces a reward function inherently independent of baselines.

\section{Experiments}
Here, we provide an empirical assessment of DPO's effectiveness in training policies from preference data. We begin by examining a controlled text-generation task to investigate the efficiency of DPO in balancing reward maximization with the minimization of KL-divergence to the reference policy, and compare this against prevalent preference learning methods such as PPO. Subsequently, we test DPO on more challenging RLHF applications, using larger model architectures on tasks like summarization and dialogue. Our findings indicate that, with minimal hyperparameter tuning, DPO generally matches or exceeds the performance of competitive methods, such as RLHF with PPO and choosing the best among $N$ sampled trajectories using a trained reward model. Prior to discussing these empirical findings, we first outline our experimental protocol; further implementation specifics are provided in Appendix~\ref{app:exp_details}.

\textbf{Tasks.} Our study investigates three distinct open-ended text generation challenges. Across all experiments, each algorithm is tasked with learning a policy using a dataset of preferences $\mathcal{D}=\bigl\{x^{(i)}, y_w^{(i)}, y_l^{(i)}\bigr\}_{i=1}^N$. For \textbf{controlled sentiment generation}, the prompt $x$ is an initial segment of a movie review from the IMDb dataset \cite{maas-EtAl:2011:ACL-HLT2011}, with the requirement that the model produce a continuation $y$ displaying positive sentiment. To facilitate a rigorous assessment, we automatically construct pairs of generations for comparison by leveraging a pre-trained sentiment classifier, ensuring that $p(\text{positive}\mid x,y_w)>p(\text{positive}\mid x,y_l)$. In the case of SFT, we adapt GPT-2-large to the task via continued training on movie reviews from the training set of the IMDb dataset, proceeding until the model stabilizes (for further information, refer to App~\ref{app:sentiment_details}). In the \textbf{summarization} task, $x$ represents a Reddit forum post, and the policy is expected to output a summary $y$ distilling the key points. We follow established conventions by employing the Reddit TL;DR summarization corpus \citep{volske-etal-2017-tl} together with human preference annotations curated by \citeauthor{stiennon2022learning}. For this setting, we leverage an SFT model further refined on human-composed forum summaries\footnote{\url{https://huggingface.co/CarperAI/openai_summarize_tldr_sft}} utilizing the TRLX \citep{leandro_von_werra_2023_7790115} implementation for RLHF. Notably, the preference annotations from \citeauthor{stiennon2022learning} were generated for outputs from a related, but distinct, SFT system. Lastly, the \textbf{single-turn dialogue} scenario presents $x$ as a user-issued prompt—ranging from technical questions to personal advice requests—and the model’s objective is to generate a response $y$ that is both useful and engaging. This is evaluated on the Anthropic Helpful and Harmless dialogue dataset \citep{bai2022training}, which contains 170,000 dialogues between people and AI assistants. Each exchange concludes with two responses generated by a sizable (but unspecified) language model and a human annotation indicating the preferred reply. In this context, as there is no publicly available SFT model, we build our SFT system by fine-tuning an existing language model solely on the human-preferred completions.

\textbf{Evaluation.} Our experimental design incorporates two primary evaluation strategies. To systematically assess the capacity of each algorithm to optimize for constrained reward maximization, we first examine their performance in a controlled sentiment generation environment. Here, we quantify performance via the trade-off frontier between realized reward and KL-divergence relative to the reference policy, leveraging direct access to the ground-truth reward metric (as operationalized by a sentiment classifier). In contrast, such a reward function is not available in authentic deployment scenarios. Thus, in practical applications, we turn to the metric of \textit{win rate} versus a baseline policy, employing GPT-4 to emulate human assessment for summary quality in summarization tasks and response helpfulness in single-turn dialogue tasks. For summarization, baseline comparisons are drawn from reference summaries found within the test set, while in dialogue, the benchmark is the human-preferred response documented in the dataset. Although some literature indicates that LMs can surpass traditional metrics as automated evaluators \citep{Chen2023ExploringTU}, we support our use of GPT-4 for evaluation with an independent human study, detailed in Sec.~\ref{sec:human-judgments}. Our analysis reveals a strong concordance between GPT-4 and human evaluations, with the agreement level between humans and GPT-4 often equaling or exceeding the agreement observed among multiple human annotators.

\textbf{Methods.} Beyond DPO, we benchmark a range of established techniques designed to train language models in alignment with human preferences. For instance, we test zero-shot prompting with \textbf{GPT-J} \citep{gpt-j} on the summarization task and employ 2-shot prompting using \textbf{Pythia-2.8B} \citep{biderman2023pythia} within the dialogue context. Moreover, we include the \textbf{SFT} model and introduce \textbf{Preferred-FT}, which utilizes supervised fine-tuning based on the selected output $y_w$—drawn from the SFT model in both sentiment and summarization, or from a general-purpose LM for dialogue. An alternative pseudo-supervised approach, \textbf{Unlikelihood}~\citep{welleck2019neural}, aims to reinforce the probability of $y_w$ while actively suppressing the likelihood of $y_l$; in our experiments, we incorporate a tunable parameter $\alpha\in[0,1]$ to weight the `unlikelihood' objective. Our evaluations also cover \textbf{PPO} \citep{schulman2017proximal}, where a reward model inferred from preference data is used, and \textbf{PPO-GT}, which directly leverages the oracle ground-truth reward available in the sentiment control setup. For sentiment-related studies, we utilize two versions of PPO-GT: a standard implementation \cite{leandro_von_werra_2023_7790115} and an enhanced variant featuring reward normalization and additional hyperparameter optimization—modifications we also apply to PPO runs utilizing learned rewards. Lastly, we analyze the \textbf{Best of $N$} strategy, which selects the highest-rewarded output among $N$ samples produced by the SFT model (or Preferred-FT in dialogue), where the ranking is determined by a reward function trained on preference data. While this approach can yield strong results by separating reward modeling from PPO optimization, its computational cost is significant for even moderate values of $N$, since generating $N$ samples per prompt is required at inference.

\subsection{How well can DPO optimize the RLHF objective?}

In standard RLHF approaches, the objective of maximizing reward subject to a KL constraint reflects a compromise between leveraging reward signals and ensuring the learned policy remains close to the reference. Thus, a fair comparison between methods should jointly consider both the reward attained and the accompanying KL divergence; pursuing higher rewards at the cost of substantially increased KL is generally suboptimal. The trade-off between reward and KL is illustrated for several methods in the sentiment domain in Figure~\ref{fig:frontier-tldr-main}. For each method, we conduct multiple experimental runs, each with distinct policy conservativeness parameters (target KL $\in\{3,6,9,12\}$ for PPO, $\beta \in \{0.05,0.1,1,5\}$, $\alpha\in\{0.05,0.1,0.5,1\}$ for unlikelihood, and various random seeds for preferred-FT), totaling 22 runs overall. At intervals of 100 training steps up to convergence, we assess each trained policy on a suite of test prompts, calculating both the mean reward according to the true reward function and the mean sequence-level KL\footnote{Specifically, the total of KL divergences computed at each time step.} relative to the reference policy $\text{KL}\left(\pi\mid \mid \pi_\text{ref}\right)$. Our results indicate that DPO consistently defines the most favorable reward-KL frontier, attaining superior rewards while preserving low KL values. This finding is particularly striking for several reasons. Primarily, while both DPO and PPO optimize an identical objective, DPO exhibits much greater efficiency; its reward-to-KL balance always outperforms PPO. Additionally, even in scenarios where PPO is provided access to true reward values (PPO-GT), DPO maintains a superior trade-off frontier.

\subsection{Can DPO scale to real preference datasets?}
\label{sec:dpo-real-datasets}
We proceed by assessing how well DPO adapts when fine-tuned on real-world datasets, focusing on tasks in summarization and single-turn dialogue. Summarization poses a challenge because standard automated metrics like ROUGE have been shown to correlate weakly with human preferences~\citep{stiennon2022learning}. Previous research has demonstrated that models fine-tuned with PPO on human-labeled preference data can generate more desirable summaries. To compare methods, we sample output completions from models using the test split of the TL;DR summarization dataset, then determine the average win rate against reference summaries in the same set. Each model's outputs are generated with temperatures ranging from 0.0 to 1.0, and the resulting win rates are depicted in Figure~\ref{fig:frontier-tldr-main} (right panel). For a fair comparison, DPO, PPO, and Preferred-FT all perform additional fine-tuning starting from the same GPT-J SFT checkpoint\footnote{\url{https://huggingface.co/CarperAI/openai_summarize_tldr_sft}}. Our findings indicate that DPO achieves a win rate close to 61\% at a temperature setting of 0.0, which surpasses PPO’s win rate of approximately 57\% at its best sampling temperature of 0.0. Additionally, DPO attains a higher peak win rate than the $N$-best baseline. It is important to highlight that we did not extensively tune the $\beta$ hyperparameter for DPO, so these outcomes might not represent its optimal performance. Furthermore, DPO displays increased robustness across a range of sampling temperatures relative to PPO, whose effectiveness diminishes to that of the original GPT-J model at higher temperatures. Notably, Preferred-FT yields negligible improvement over the base SFT model. In Section~\ref{sec:human-judgments}, we extend this analysis through direct human preference judgments, where completions from DPO sampled at temperature 0.25 were favored 58\% of the time when compared to PPO outputs at the same temperature.

For single-turn dialogue evaluation, we assess the various approaches on a subset of the Anthropic HH dataset test split \citep{bai2022training} that features a single round of human-assistant exchange. In our GPT-4 assessments, we utilize the test set’s preferred completions as references to determine the win rate for each method. Since there is no established SFT model tailored to this task, our experiments begin with a pre-trained Pythia-2.8B model. We employ Preferred-FT to fine-tune a reference model on the selected completions, ensuring that generated completions remain within the distribution of the model, before proceeding with DPO training. Additional comparisons include the strongest result from 128 Preferred-FT completions (our experiments indicate that the Best of $N$ strategy reaches a plateau at 128 completions in this context; refer to Appendix Figure~\ref{fig:best-of-n}) and a 2-shot prompt variant of the base Pythia-2.8B model. Results indicate that DPO matches or surpasses other methods at their respective optimal temperatures.

We further investigate an RLHF model trained via PPO on the Anthropic HH dataset \footnote{\url{https://huggingface.co/reciprocate/ppo_hh_pythia-6B}}, sourced from a recognized implementation \footnote{\url{https://github.com/CarperAI/trlx/tree/main/examples/hh}}. Despite this, we are unable to identify a prompt or sampling temperature that leads to better outcomes than those produced by the unmodified Pythia-2.8B. Drawing from observations in TL;DR and acknowledging that both Best of 128 and PPO directly optimize the same reward signal, we treat the Best of 128 metric as an approximate indicator of PPO's performance. In summary, DPO stands out as the only method that is both computationally efficient and capable of exceeding the quality of the reference completions from the Anthropic HH dataset, while also delivering comparable or improved performance relative to the more resource-intensive Best of 128 baseline. Lastly, as illustrated in Figure~\ref{fig:dialogue-main}, DPO attains peak performance in a relatively short amount of training time.

\subsection{Generalization to a new input distribution}

To further explore how PPO and DPO policies handle distributional changes, we assess the models trained on Reddit TL;DR summarization when applied to a new data domain—specifically, news articles from the CNN/DailyMail test set \citep{nallapati-etal-2016-abstractive}. For this evaluation, we used the optimal sampling temperatures identified from the TL;DR task (0 and 0.25). The comparative results are summarized in Table~\ref{tab:ood}. GPT-4 win rates were calculated using the same evaluation setup as for the Reddit TL;DR domain, except that the phrase ``forum post'' was swapped for ``news article'' in the GPT-4 (C) prompt. In this cross-domain scenario, DPO policies consistently outperformed their PPO counterparts by a substantial margin. These findings offer preliminary support that DPO can generalize to novel distributions comparably to PPO—even though DPO training omits the extra, unlabeled Reddit TL;DR prompts utilized by PPO.

\subsection{Validating GPT-4 judgments with human judgments}
\label{sec:human-judgments}

A human evaluation study was carried out to assess the reliability of GPT-4's comparative judgments, leveraging the outputs from the TL;DR summarization task and employing two distinct GPT-4 prompts. The \textbf{GPT-4 (S)} (simple) prompt asks which summary most effectively captures the main information from the post. Meanwhile, the \textbf{GPT-4 (C)} (concise) prompt solicits judgments not only on informativeness but also conciseness, motivated by observations that GPT-4's preferences using the \textbf{GPT-4 (S)} prompt tend to favor summaries that are longer and more repetitive compared to human raters. Full details of these prompts are available in Appendix~\ref{app:prompts}. To provide a range of sample qualities, we compare three approaches: the best-performing (DPO, temp. 0.25), the lowest-performing (PPO, temp. 1.0), and an intermediate method (SFT, temp. 0.25), each against the PPO model using greedy sampling (its best setting). Our analysis reveals that for both prompts, GPT-4's preference alignment with human raters is roughly on par with inter-human agreement rates, indicating GPT-4's viability as a stand-in for human evaluators (noting that, because of limited numbers of human raters, multiple independent human judgments were obtained only for DPO and PPO-1 cases). In general, the \textbf{GPT-4 (C)} prompt produced win rates most aligned with human judgments, which justifies its selection as the primary metric in Section~\ref{sec:dpo-real-datasets}. For further information regarding the human evaluation protocol, including details about the rater interface and participant roster, consult Appendix~\ref{app:human-study}.

\section{Discussion}
Preference-based learning constitutes an effective and scalable method for developing competent and aligned language models. In this work, we introduced DPO, a straightforward approach that facilitates training language models from preference data without relying on reinforcement learning. Unlike prior approaches that adapt preference-based learning into a conventional RL framework to take advantage of existing RL techniques, DPO instead establishes a direct correspondence between language model policies and reward representations. This approach permits optimizing language models to align with human preferences using a simple cross-entropy loss, bypassing both reinforcement learning and the limitations associated with generality. DPO demonstrates performance comparable to, and occasionally surpassing, traditional RLHF methods such as those utilizing PPO, all with minimal hyperparameter adjustment. Consequently, DPO significantly streamlines the process of preference-based language model training.

\textbf{Limitations \& Future Work.} Our findings open up several avenues for further exploration. A key question is how DPO-trained policies perform on out-of-distribution data compared to those trained with explicit reward modeling; preliminary evidence suggests that DPO and PPO-based methods achieve similar generalization, though more systematic analysis is necessary. Furthermore, it is worth investigating whether self-labeling strategies applied during DPO training can leverage unannotated prompts as effectively. Another area of inquiry concerns the manifestations of reward over-optimization within the DPO framework, including whether the modest performance decline observed in Figure~\ref{fig:dialogue-main}-right results from such phenomena. Although our current experiments span models up to 6B parameters, scaling DPO to accommodate models several orders of magnitude larger represents a promising research direction. Regarding assessment protocols, our experiments indicate that GPT-4's win rate metrics are sensitive to prompt variation; future studies should aim to optimize prompt design for robust automated evaluation. Beyond these considerations, DPO offers broad applicability not only to preference-based language model training but also to the development of generative models across diverse modalities.